\subsection*{V9. Spin waves|Magnons|Dispersion Relation for Spin wave}
\addcontentsline{toc}{subsection}{9. Spin waves|Magnons|Dispersion Relation for Spin waves}
\begin{enumerate}
    \item Analogous to lattice waves.
    \item Generated due to the exchange forces.
    \item \hl{Collective motion of magnetic moment in magnetically ordered materials like ferromagnetic materials}.
    \item Due to some external disturbance, these magnetic moment goes to some motion which give rise to spin-waves.
    \item Can be also called as the wave of a quantized energy that propagates through a substance as a result of magnetic field shifts within an atom in response to an outside stimulus (alternating field).
    \item Quantization of spin waves gives us \hl{\textbf{Magnons}}.
    \item \textbf{Magnons:} 
    \begin{enumerate}[(i.)]
        \item Collective excitations of electrons' spin structures in a crystal lattice, crystal should be magnetic in nature, where the order state is observed, i.e., ferromagnetic or antiferromagnetic and also ferrimagnetic substances.
        \item Magnons are a quasi particle.
        \item Quasi particle means the particle which can only exists when interaction is taking place, that means inside the interacting many particle systems.
        \item They are Bosons and obeys Bose-Einstein theory.
        \item Spin = 1 (integral spin)
        \item First introduced by Felix Bloch in 1930, used to introduce to explain the reduction of spontaneous magnetization in ferromagnets.
        \item $\mathrm{T=0}$ K, lowest energy state with parallel atomic spins.
        \item As the temperature is increased, there is increase in the internal energy (randomness in the system) of the system, and decrement in the spontaneous magnetization which as a result \textbf{magnon} is created.
        \item Each magnon will reduce total spin by 1 $\mathrm{\hbar}$ and magnetization is reduced by $\mathrm{\gamma\hbar}$, where $\mathrm{\gamma}$ is the gyromagnetic ratio\footnote{Ratio of magnetic moment to the angular momentum.}.
        \item Experimentally detected in 1957 by Bertram Brockhouse by inelastic neutron scattering in ferrite.
        \item Detected in Ferro-, Ferri-, and Anti-ferromagnets.
    \end{enumerate}
    \item Dispersion Relation for Spin waves:
    \begin{enumerate}[(i.)]
        \item Ferromagnetic ground state is ordered one, hence magnon is not created. 
    \end{enumerate}
    \item Quantized lattice waves gives us \hl{\textbf{Phonons}}.     
\end{enumerate}